% !PRINTED 535
\weberpage{535}
%<<CONT
setzen (wo jetzt $A$, $B$ natürlich nicht zu verwechseln sind mit den
Koeffizienten der quadratischen Form), so ist $B$ eine ganze algebraische
Zahl, und nach § 74, (14) ist $A$ die Wurzel einer algebraischen
Gleichung, deren Koeffizienten ganze algebraische Zahlen sind. Es ist
also $A$ ebenfalls eine ganze algebraische Zahl.

Da aber die beiden Quotienten $u^r:v^r$ und $v^r:u^r$ die Wurzeln der
quadratischen Gleichung
\[
x^2 - Ax + 1 = 0
\]
sind, so folgt, daß
\[
\frac{u}{v} \quad \text{und} \quad \frac{v}{u}
\]
ganze Zahlen, und da sie zueinander reziprok sind, \gesperrt{Einheiten sind.}

\gesperrt{Es sind also $u$ und $v$ assoziierte Zahlen.}

Da man nun nach § 118 durch wiederholte Kompositionen mit Klassen $l$
(durch die Primzahlen darstellbar sind) von jeder Klasse $k$ zu jeder
anderen Klasse $k'$ derselben Diskriminante gelangen kann, und da zwei
mit einer dritten assoziierten Zahl untereinander assoziiert sind, so
haben wir den Satz:

\gesperrt{Setzt man für $\omega$ die $h$ Wurzeln der
Formen eines Systems von Repräsentanten der Diskriminante $-4m$, so sind
die $h$ Zahlen $f(\omega)$ untereinander assoziiert}; und daraus nach (1)
unmittelbar den merkwürdigen Satz, der sich leicht an allen Beispielen
bestätigen läßt:

\gesperrt{$f(\omega):2^\tau$ ist eine ganze Zahl, und
zwar eine Einheit.}

1. Die Bestimmung der Exponenten $\tau$ ist durch elementare
Hilfsmittel möglich, wenn $m \equiv 1$ oder $\equiv 2\ (\mathrm{mod}\ 4)$
oder $m \equiv 3\ (\mathrm{mod}\ 8)$.

Machen wir in der Gleichung mit ungeraden äußeren Koeffizienten
\weq{3}{A\omega^2 + 2B\omega + C = 0}
mit der Diskriminante $4(B^2 - AC) = -4m$ die Substitution
\weq{4}{\omega = \frac{\omega'-1}{\omega'+1}, \quad \omega' = -\frac{\omega+1}{\omega-1},}
so erhalten wir die Gleichung
\weq{5}{\frac{A+2B+C}{2}\,\omega'^2 - (A-C)\,\omega' + \frac{A-2B+C}{2} = 0,}
%CONT>>
